\chapter{Fragen}


\begin{enumerate}
	\item 
	In Sundberg Section 3: Hier wird eine Methode vorgestellt, wie man ohne Simulation die non coverage Wahrscheinlichkeit berechnen kann.
	Das ist die Grundlage für die Figures 2-5. Wenn man das Konfidenzintervall $CI = [c_l \hat \theta, c_u \hat \theta]$ betrachtet (C1, C2, C3, C4, C6, C7), dann ist die non coverage W'keit gegeben durch.
	\[
		\mathbb{P} \left( \text{non coverage} \right)
		=
		1 - \mathbb{P} \left( \text{coverage} \right)
		=
		1 - \mathbb{P} \left( c_l \hat \theta \leq \theta \leq c_u \hat \theta \right)
		=
		1 - 
			\mathbb{P} \left( \theta \leq c_u \hat \theta \right)  
			+ 
			\mathbb{P} \left( \theta \leq c_l \hat \theta \right) 
	\]
	Den Term $\mathbb{P} \left( \theta \leq c_l \hat \theta \right)$, kann man mit der ersten Equation in Section 3 von Sundberg mit $t = \hat \theta / c_l$ berechnen.
	\[
		\mathbb{P} \left( \theta < c_l \hat \theta  \right) 
		=
		\mathbb{P} \left( \theta / c_l < \hat \theta  \right)
		=
		1 - \mathbb{P} \left( \theta / c_l > \hat \theta  \right)
	\]
	
	Jedoch lebt diese Equation davon, dass man das Gesetz der totalen Wahrscheinlichkeit auf $N$ anwendet, doch nun ist $t$ auch abhängig von $N$, weil $c_l$ von $N$ abhängig ist.
	Das macht die Berechnung sehr viel schwerer.
	
	In Sundberg benutzt man $F(. | N)$, nämlich
	\[
		\mathbb{P} \left( \hat \theta < t \right)
		=
		\mathbb{E}_N
		\left[
			F(Nt - (n-N)C | N)
		\right]
	\]
	doch in unserem Fall, hängt $t$ auch von $N$ ab.
	\textcolor{red}{($\star$) Idee: einfach beim Erwartungswert die Summe nehmen und dann auch $t$ in Abhängigkeit von $N$ darstellen.}
	
	\item Wie kann ich die Idee von Section 3 in Sundberg auf C5 anwenden? C5 hat keine Form wie $CI = [c_l \hat \theta, c_u \hat \theta]$. Wie kann ich dann trotzdem die non coverage W'keit ohne Simulation bestimmen?
	
	\item 
	Welcher wahre Parameter für $\theta$ wird in Sundberg Figures 2-5 verwendet? Ist das überhaupt notwendig zu wissen, um die Plots zu erstellen.
	
	\item 
	Für Sundberg C2 gilt $c_l > c_u$ für $N \in \{1, 2\}, \alpha = 0.1$ und damit ist das Konfidenzintervall nicht valid
\end{enumerate}